\section{Conclusiones}
La construcción de comunidades científicas mediante prácticas abiertas requiere conectar tareas realizables, acompañamiento, revisión y conservación. CoAfina, Salud para todos y las experiencias educativas con datos de física permiten examinar esas conexiones desde productos y escalas temporales diferentes. Su valor conjunto no reside en sumar participantes, sino en mostrar decisiones organizativas que pueden explicitarse y adaptarse.

El marco propuesto distingue participación, producción, reutilización y continuidad para evitar que un indicador sustituya a los demás. Para instituciones venezolanas ofrece un punto de partida: seleccionar una pregunta abordable, asignar responsabilidades, preparar materiales reutilizables y documentar qué sucede durante y después de la actividad. La siguiente etapa consiste en aplicar definiciones compartidas y reunir evidencia prospectiva que permita evaluar tanto las oportunidades creadas como sus límites.

\section*{Procedencia y estado del manuscrito}
Este borrador desarrolla el resumen ICHEP 2026 número 668, aceptado como charla y posteriormente retirado según comunicación de los autores por imposibilidad de asistencia. El arbitraje del resumen no constituye arbitraje de este texto. Se citan los antecedentes publicados y no se reivindican sus resultados como inéditos. La validación de autoría, afiliaciones, financiación, conflictos de interés y direcciones de correspondencia permanece pendiente de revisión de los autores. El registro de evidencia se incluye en el repositorio acompañante.
